% !TeX root = ../papers/nc-p3.tex
\section{三十節点を印した四次式の成分と中心写像}
\label{sec:thirty-node-center-map}
\begin{theorem}\label{thm:thirty-node-center-map}
十五次元成分上の一般窓から射影中心を取る写像は、三十節点四次三次元多様体の
一つの既約成分へ生成的有限かつ支配的で、一般点で不分岐である。
一般窓の局所モジュライは、有限の印を付ければ中心の三十節点を保つ変形の局所モジュライ
とエタール局所で同型である。中心は $\mathbb P^4$ 内の順序なしの三十点だけからも一意に復元できる。
\end{theorem}
\begin{proof}
$S,J$ を前節と同じとし、$\ell=a+2b+3c+4d+5z$ と置く。
同じ法 $1009^5$ の中心で $J_7+\ell^4 S_3$ の行列を作ると、単元ピボットは
付値 $0,1,2$ にそれぞれ $298,0,32$ 個ある。
従って一般点で $J_7+\ell^4 S_3=S_7$ である。
\path{computations/verify_thirty_node_center_map.py} が入力四次式の構成も含めて検算する。
特に $S_7=J_7+\ell S_6$ であり、両辺に $S_1$ を掛け続ければ
$n\ge7$ で $S_n=J_n+\ell^{n-3}S_3$ となる。
これから三十節点集合 $D$ のどの点でも $\ell\ne0$ である。
$n$ を大きくすれば $S_n\to H^0(D,\mathcal O_D(n))$ は全射で、$J_n$ はその核に入る。
従って $S_3$ の評価は全射である。特に五偏微分は三十点を通る三次式を全て張る。
実際 $\dim I_D(3)=35-30=5$ で、偏微分が従属なら四次式は錐となり、
その頂点は通常二重点でなくなる。従って $I_D(2)=0$ でもある。
二次式があればその五つの一次倍が全偏微分を張り、特異集合が曲面以上の次元を持つからである。

四次式 $G$ が同じ三十点で特異なら $dG=A\,dF$ となる定数行列 $A$ がある。
混合偏微分の対称性により $A$ の転置は $F$ のHessianに関して自己随伴である。
異なる一般固有空間は任意の変数値で直交するので、二固有値があれば $F$ は
互いに異なる変数群の四次式の和になる。その特異点が二群以上に非零成分を持てば
成分を別々に拡縮して正次元の特異集合を得る。一群にしか成分がなければ、他群の全方向が
アフィン接チャートのHessianの核に残る。いずれも三十節点に反する。
従って固有値は一つである。非零の冪零部 $N$ の指数が $s\ge2$ なら、
$v=N^{s-1}w\ne0$ を取り、自己随伴性と $N^{2s-2}=0$ より $D_v^2F=0$ を得る。
$v$ 方向の変数について一次以下の斉次四次式は $[v]$ で重複度が少なくとも3である。
これも矛盾するので $A$ はスカラー、$G$ は $F$ のスカラー倍である。
中心の射影自己同型も有限である。三十点を一つずつ固定する非自明な射影変換の
固有空間の合併は二超平面の和に含まれる。一固有値で非半単純なら一超平面で足りる。
これは $I_D(2)=0$ に反する。

次に中心を固定した窓の一次変形を考える。$X$ 上の折り畳み代数層を $\mathcal B$ とする。
$k[\epsilon]/(\epsilon^2)$ 上の中心自明化を一つ指定する。
各閉点で中心局所環を完備化すれば、滑らかな所では $\mathcal B$ は行列環、
節点では $A=\operatorname{End}_R(M)$、$M=R^a\oplus I^b$ である。
後者で $M$ は $R$ 上の自己第一拡大を持たない。$R$ を第一引数とする場合は自由性、
第二引数とする場合はGorenstein環上での極大Cohen--Macaulay性による。
$I$ の自己拡大の中項は反射的である。閉点を除けば $I$ は直線束で
$H^1(U,\mathcal O)=0$（$R$ の深さ3）なので拡大は分裂する。
分裂射は反射性により閉点を越えて延びる。これで $\operatorname{Ext}^1_R(M,M)=0$ が従う。

$R$ 加群 $R\subset M$ を一つ選ぶ原始冪等元 $p$ を一次変形 $A'$ に持ち上げる。
$A'p$ は $k[\epsilon]/(\epsilon^2)$ 平坦で、剰余が $M$ なので、先の消滅により
$R[\epsilon]/(\epsilon^2)$ 加群として $M[\epsilon]/(\epsilon^2)$ と同型である。
左乗法による $A'\to\operatorname{End}(A'p)$ は剰余で標準同型で、両辺が
双対数上平坦なので同型である。滑らかな点でも $M$ が自由として同じ証明になる。
従って代数層の変形は局所自明である。完備化から局所環へ降ろすには、まず
$R$ 加群の自己拡大類の消滅を忠実平坦性で降ろす。次に加群分裂を固定すれば、
積の一次補正を自明にする条件は有限加群の準同型についての線型方程式であり、
その余核の類の消滅も同様に降りる。

中心上の導分の層は $\mathcal B/\mathcal O_X$ である。実際一般行列代数上の導分は内的で、
被約trace零の実現元は一意である。余次元1ではAzumayaであり、内的導分の表示が
エタール局所の行列代数から降りる。従ってそのtrace零の元は全ての高さ1で正則で、
$\mathcal B$ の反射性により全体でも正則である。
$\mathcal B$ の各隅には中間コホモロジーがない（中心上の次数付き極大Cohen--Macaulay性）。
また $H^2(X,\mathcal O_X)=0$ なので、
$H^1(X,\mathcal B/\mathcal O_X)=0$。
従って局所自明化の差の一次導分を貼り合わせれば変形は大域でも自明である。
四つの順序付き冪等元の二つの持ち上げは大域切断環内で $1+\epsilon b$ の共役で移り合う
（対応する新旧冪等元の積の和がその共役元になる）。
四連続射影の順方向のHomとその積はこの代数層の対応する隅の全切断である。
特に最初の三文字と二つの隣接関係を回収でき、系\ref{cor:two-window-extension}で
全窓が復元される。この復元は双対数上でも単元小行列によって同じである。
これで中心を取る微分の核は文字基底変更だけである。

最後に番号付き三十節点と四次式の接続多様体を考える。
各節点を保つ一次条件は四次変分がその点で消えることであり、点の移動は
可逆な四次元Hessianから一意に決まる。実際各点の局所チャートで値と四偏微分を零にする
五式を使えばよい。三次式の評価が全射なので四次式の評価も全射であり、
この接続多様体は射影四次式空間の次元69から30を引いた滑らかな次元39である。
番号の付け替えは有限で、余分の特異点を持たない条件はここで開である
（印の近傍ではHessianで特異点を一意に解け、補集合には固有性を適用する）。
射影五座標の変更は有限安定群で次元24なので、商の次元は15となる。
窓側は滑らかな次元15の有理商を持つ。両側で必要なら三十点を番号付けし
局所座標枠を取る。滑らかな同次元空間の間の単射微分から、得た中心写像は
一般点でエタールであり、その像はこの成分で稠密である。有限型性により一般繊維は有限である。
\end{proof}
